% =============================================================================
%  Capítulo 2 -- Marco teórico
%  -------------------------------------------------------------------------
%  NOTA PARA AXEL:
%  - Este capítulo está redactado con contenido real, pero deliberadamente
%    sintético: cada sección es un primer borrador sólido que debes expandir
%    con detalle de las fuentes primarias.
%  - Las marcas «[VERIFICAR: ...]» señalan afirmaciones cuya formulación
%    exacta conviene cotejar contra el artículo citado antes de la entrega.
%  - Las marcas «[FUENTE MX: ...]» señalan dónde DEBES añadir bibliografía
%    mexicana específica (esto cubre tu objetivo específico O1). No inventé
%    referencias para México: no existen aún en referencias.bib.
% =============================================================================
\chapter{Marco teórico}
\label{cap:marco}

Este capítulo construye el andamiaje conceptual de la tesis articulando cuatro
cuerpos de literatura que rara vez se leen en conjunto. La
Sección~\ref{sec:sesgo-algoritmico} sitúa la noción de sesgo en sistemas de
inteligencia artificial y precisa el sentido en que se usará a lo largo del
trabajo. La Sección~\ref{sec:pln} introduce los elementos técnicos de la
clasificación automática de texto necesarios para interpretar los experimentos
posteriores. La Sección~\ref{sec:sociolectos} recupera, desde la
sociolingüística, los conceptos de variación sociolectal y discriminación
lingüística. Finalmente, la Sección~\ref{sec:clasismo} aterriza estas
discusiones en el contexto mexicano, que es el escenario empírico de la
investigación.

\section{Sesgo algorítmico en sistemas de inteligencia artificial}
\label{sec:sesgo-algoritmico}

\subsection{Tres sentidos de ``sesgo''}
\label{sec:tres-sentidos}

El término \emph{sesgo} es notoriamente polisémico, y buena parte de la
confusión en los debates públicos sobre inteligencia artificial proviene de no
distinguir sus acepciones. Conviene separar al menos tres.

En primer lugar, el \emph{sesgo estadístico} designa la diferencia sistemática
entre el valor esperado de un estimador y el parámetro que pretende estimar. Es
una noción técnica, neutra desde el punto de vista moral, y constitutiva de
todo aprendizaje: el llamado \emph{compromiso sesgo--varianza} muestra que
cierto sesgo es incluso deseable para controlar la varianza de un modelo. En
segundo lugar, el \emph{sesgo cognitivo} se refiere a las desviaciones
sistemáticas del juicio humano respecto de la racionalidad normativa, estudiadas
por la psicología. En tercer lugar, el \emph{sesgo social} o
\emph{discriminatorio} designa el trato sistemáticamente desigual de grupos
sociales por parte de un sistema, en formas que reproducen o amplifican
desigualdades preexistentes \cite{Mehrabi2021Survey}.

Esta tesis se ocupa exclusivamente del tercer sentido, cuya relevancia para el
\PLN fue señalada tempranamente por Hovy y Spruit \cite{HovySpruit2016} al
argumentar que las tecnologías del lenguaje tienen un impacto social propio
---sobre la exclusión, la sobregeneralización y la exposición de grupos---
que la comunidad técnica había tendido a ignorar. La distinción entre los tres
sentidos no es ociosa: un sistema puede estar estadísticamente bien calibrado
y, sin embargo, ser socialmente sesgado, porque los datos sobre los que aprende
ya codifican desigualdades del mundo. La cuestión, por tanto, no es meramente técnica sino
\emph{sociotécnica}: requiere examinar conjuntamente el comportamiento del
modelo y el contexto social en el que se inserta.

\subsection{Definición operativa}
\label{sec:definicion-sesgo}

Para que el problema sea evaluable empíricamente se adopta la siguiente
caracterización operativa.

\begin{definicion}[Sesgo social en un sistema de clasificación]
\label{def:sesgo}
Sea $f : \mathcal{X} \to \mathcal{Y}$ un clasificador y sea $A$ un atributo
protegido o socialmente sensible definido sobre la población de entradas. Se
dice que $f$ exhibe \emph{sesgo social} respecto a $A$ si existen valores
$a, a' \in A$ y una métrica de comportamiento $\mu$ tales que
\[
  \mu\bigl(f \mid A = a\bigr) \;\neq\; \mu\bigl(f \mid A = a'\bigr),
\]
y dicha diferencia no se justifica por una variación correspondiente en la
variable objetivo $\mathcal{Y}$ atribuible legítimamente a $A$.
\end{definicion}

La cláusula final ---``no se justifica por una variación legítima''--- es la que
carga con todo el peso normativo y, a la vez, la más difícil de operacionalizar.
En el caso de esta tesis, el atributo $A$ es el \emph{registro lingüístico} de
un texto, y la justificación se vuelve especialmente tenue: si dos textos
comunican el mismo contenido sustantivo y difieren solo en su estilo, no existe
ninguna variación legítima en $\mathcal{Y}$ que pueda explicar un trato
diferenciado. Es precisamente este caso límite el que el \emph{Índice de
Neutralidad Sociolectal} (\INS), definido en el Capítulo~\ref{cap:ins}, busca
capturar.

\subsection{Fuentes de sesgo en sistemas de procesamiento de lenguaje natural}
\label{sec:fuentes-sesgo}

Hovy y Prabhumoye \cite{HovyPrabhumoye2021} ofrecen una taxonomía influyente que
identifica cinco puntos del ciclo de vida de un sistema de \PLN en los que puede
introducirse sesgo social. Conviene revisarlos porque ofrecen un mapa de dónde
podría originarse el sesgo sociolectal que la tesis investiga.

\subsubsection{Sesgo en los datos}

Los modelos de \PLN aprenden de corpus que son muestras del lenguaje producido
en condiciones sociales concretas, no espejos neutrales de la lengua. El sesgo
de los datos surge cuando ciertos grupos, registros o variedades están
sobre- o subrepresentados, o cuando el texto disponible refleja desigualdades
históricas. A medida que los corpus de entrenamiento crecen hasta volúmenes que
ningún equipo puede auditar exhaustivamente, este problema se agrava: Bender
y colaboradores \cite{Bender2021Parrots} advierten que el tamaño no garantiza
diversidad y que los corpus masivos tienden a sobrerrepresentar las variedades
hegemónicas y los puntos de vista dominantes. [VERIFICAR: matiz exacto del
argumento sobre tamaño vs.\ diversidad.]

\subsubsection{Sesgo en las anotaciones}

Gran parte de los sistemas de clasificación se entrena con etiquetas asignadas
por anotadores humanos, cuyas intuiciones no son neutrales. La composición
demográfica del panel de anotación, las guías de etiquetado y los juicios
implícitos sobre qué cuenta como ``correcto'' o ``ofensivo'' se transfieren al
modelo. Un ejemplo documentado y directamente pertinente para esta tesis es el
de Sap y colaboradores \cite{Sap2019Risk}, quienes mostraron que los anotadores
tendían a clasificar como ofensivo el inglés afroamericano vernáculo con mayor
frecuencia que el inglés estándar, sesgo que los clasificadores entrenados sobre
esas anotaciones aprendían y reproducían. El caso ilustra el mecanismo central
del presente trabajo: una variedad lingüística asociada a un grupo social
recibe un trato sistemáticamente distinto que no responde al contenido del
mensaje.

\subsubsection{Sesgo en las representaciones lingüísticas}

Antes de ser clasificado, el texto se transforma en representaciones numéricas.
Estas representaciones no son inertes: codifican regularidades estadísticas del
corpus, incluidas las asociaciones estereotipadas. Caliskan y colaboradores
\cite{Caliskan2017Semantics} demostraron que los \emph{embeddings} de palabras
entrenados sobre texto de la web reproducen sesgos humanos medibles, análogos a
los que captan las pruebas de asociación implícita en psicología. Bolukbasi y
colaboradores \cite{Bolukbasi2016Man} ilustraron el fenómeno con analogías de
género estereotipadas y propusieron técnicas tempranas de mitigación. Estos
trabajos establecen que el sesgo puede residir en la capa de representación,
con independencia del clasificador final.

\subsubsection{Sesgo en el diseño del modelo}

La arquitectura, la función objetivo y los hiperparámetros de un modelo también
introducen sesgo, en particular porque los modelos tienden a \emph{amplificar}
las regularidades del entrenamiento más que a reproducirlas fielmente. Un
patrón presente en una minoría de los datos puede volverse la respuesta por
defecto del modelo si reduce la pérdida agregada. Para esta tesis, la
implicación es metodológica: el sesgo sociolectal observado en un clasificador
no puede atribuirse sin más a sus datos de entrenamiento, pues la magnitud del
efecto depende también de la familia de modelo; de ahí que el diseño
experimental del Capítulo~\ref{cap:metodologia} contemple modelos de
complejidad creciente.

\subsubsection{Sesgo en el diseño de la investigación}

La fuente más sutil es la elección misma de qué problemas se estudian, sobre qué
lenguas y para qué poblaciones. La concentración del \PLN en el inglés y en
contextos del Norte global determina qué fenómenos se vuelven visibles y cuáles
permanecen sin estudiar. Esta tesis es, en parte, una respuesta a ese sesgo de
agenda: investiga la variación sociolectal del español mexicano, un objeto
sistemáticamente ausente de los \emph{benchmarks} de equidad disponibles.

\subsection{Del sesgo a la equidad: una nota crítica}
\label{sec:critica-bias}

Blodgett y colaboradores \cite{Blodgett2020Survey} señalan que buena parte de la
literatura sobre ``sesgo'' en \PLN adolece de imprecisión conceptual: se invoca
el término sin especificar qué comportamiento es perjudicial, para quién y por
qué. Recomiendan anclar todo estudio de sesgo en un daño concreto y en su
relación con las desigualdades sociales que lo motivan. Esta tesis adopta esa
recomendación: el daño concreto es el trato diferenciado de hablantes según su
registro en procesos de consecuencias materiales ---como el filtrado de
currículums--- y la desigualdad de fondo es el clasismo lingüístico discutido
en la Sección~\ref{sec:clasismo}.

\section{Procesamiento de lenguaje natural y clasificación automática de texto}
\label{sec:pln}

\subsection{La tarea de clasificación de texto}
\label{sec:tarea-clasificacion}

La clasificación de texto es la tarea de asignar a un documento $x \in
\mathcal{X}$ una etiqueta $y \in \mathcal{Y}$ de un conjunto finito y
predefinido. Formalmente, se busca aprender una función $f : \mathcal{X} \to
\mathcal{Y}$ a partir de un conjunto de ejemplos etiquetados. La canalización
(\emph{pipeline}) típica comprende cuatro etapas: (i)~\emph{preprocesamiento}
del texto crudo; (ii)~\emph{representación} del documento como un vector o una
secuencia de vectores; (iii)~\emph{modelado}, donde un algoritmo de aprendizaje
estima una función de decisión; y (iv)~\emph{decisión}, donde se emite la
etiqueta y, opcionalmente, una distribución de probabilidad sobre las clases.
El \INS, definido más adelante, opera sobre las salidas de las etapas (iii) y
(iv): la decisión y la probabilidad posterior.

\subsection{Representaciones del texto}
\label{sec:representaciones}

La forma de representar el texto determina qué información queda disponible para
el clasificador. Tres familias resultan relevantes para esta tesis.

Las representaciones \emph{de bolsa de palabras} y su ponderación
\emph{TF--IDF} describen un documento por las frecuencias de los términos que
contiene, ignorando el orden. Son interpretables y computacionalmente baratas,
pero ciegas al contexto: dos textos con el mismo contenido y distinto registro
pueden recibir representaciones muy distintas solo por su léxico, lo que las
hace un punto de partida útil para estudiar el sesgo sociolectal.

Los \emph{embeddings} estáticos asignan a cada palabra un vector fijo aprendido
de grandes corpus, de modo que palabras de uso similar quedan próximas en el
espacio vectorial. Como se discutió en la Sección~\ref{sec:fuentes-sesgo}, estos
vectores codifican asociaciones sociales del corpus de origen.

Los \emph{embeddings} contextuales, popularizados por arquitecturas de tipo
\emph{transformer} como BERT \cite{Devlin2019BERT}, asignan a cada aparición de
una palabra un vector que depende de su contexto, lo que mejora notablemente el
desempeño en la mayoría de las tareas. Para el español existen modelos
preentrenados específicos, como BETO \cite{CanetePilar2020BETO}, que serán
candidatos naturales para la evaluación experimental del
Capítulo~\ref{cap:resultados}.

\subsection{Familias de modelos relevantes}
\label{sec:familias-modelos}

Para los fines de la tesis conviene distinguir tres familias de clasificadores,
con grados crecientes de complejidad y opacidad. Los \emph{modelos lineales}
(regresión logística, máquinas de vectores de soporte) sobre representaciones
TF--IDF ofrecen interpretabilidad y un comportamiento predecible. Los
\emph{modelos basados en \emph{transformers} preentrenados} y ajustados a la
tarea, como BETO \cite{CanetePilar2020BETO}, alcanzan mayor desempeño a costa de
transparencia. Por último, los \emph{modelos de lenguaje grandes} empleados como
clasificadores en modo \emph{zero-shot} o \emph{few-shot} permiten clasificar sin
entrenamiento específico, pero su comportamiento ante variaciones de registro
es poco estudiado y especialmente difícil de auditar \cite{Bender2021Parrots}.
La inclusión de al menos un representante de cada familia permitirá observar si
el sesgo sociolectal varía con la complejidad del modelo.

\subsection{Métricas estándar de evaluación y sus límites}
\label{sec:metricas}

El desempeño de un clasificador suele resumirse con la exactitud, la precisión,
la exhaustividad (\emph{recall}) y su media armónica, la medida F1, además de
métricas de calibración cuando interesan las probabilidades. Estas métricas son
\emph{agregadas}: describen el comportamiento promedio sobre un conjunto de
prueba. Su limitación, central para esta tesis, es que un modelo puede exhibir
una F1 alta y, simultáneamente, comportarse de manera sistemáticamente distinta
ante subgrupos o registros, porque los errores se compensan en el agregado
\cite{Blodgett2020Survey, Mehrabi2021Survey}. De ahí la necesidad de métricas
\emph{desagregadas} y de indicadores específicos de consistencia como el \INS,
que no miden cuánto acierta el sistema sino cuán establemente se comporta ante
variaciones que no deberían importar.

\section{Discriminación lingüística y sociolectos}
\label{sec:sociolectos}

\subsection{Variación sociolectal: conceptos básicos}
\label{sec:variacion}

La sociolingüística ha documentado de manera consistente que el lenguaje varía
de forma sistemática según factores sociales. Conviene precisar tres términos
que se usarán con cuidado en lo que sigue. Un \emph{dialecto} es una variedad
asociada principalmente a un grupo geográfico. Un \emph{sociolecto} es una
variedad asociada a un grupo social ---definido por clase, nivel educativo,
ocupación u otros factores estructurales. Un \emph{registro} es la variación que
un mismo hablante despliega según la situación comunicativa, desde lo más formal
a lo más coloquial. Estas dimensiones se entrecruzan: un rasgo puede ser a la
vez marcador sociolectal y de registro.

Lo crucial para esta tesis es que la variación sociolectal es, en lo esencial,
variación de \emph{forma} y no de \emph{contenido}: dos hablantes de
sociolectos distintos pueden comunicar exactamente la misma información. Cuando
un sistema responde de manera diferente a ambos, la diferencia no puede
atribuirse al mensaje, sino al cómo se dice ---es decir, a un marcador social.

\subsection{La discriminación lingüística como fenómeno social}
\label{sec:discriminacion-ling}

La \emph{discriminación lingüística} ocurre cuando determinados registros,
dialectos o sociolectos se asocian con juicios negativos sobre la capacidad, la
inteligencia o la credibilidad de quien los emplea, de modo que la forma del
habla se convierte en vía para reproducir jerarquías sociales preexistentes
\cite{Weissler2022}. Drożdżowicz \cite{Drozdzowicz2024} subraya que el fenómeno
es más complejo de lo que sugiere su nombre, pues entrelaza juicios estéticos,
epistémicos y morales que resultan difíciles de separar analíticamente.
[VERIFICAR: caracterización precisa de la tesis de Drożdżowicz; cotejar contra
el artículo antes de la entrega.] Lo que importa retener es que la
discriminación lingüística no es una reacción a deficiencias comunicativas
reales ---el mensaje se entiende perfectamente--- sino una sanción social
adherida a marcadores de grupo.

[FUENTE MX: aquí conviene introducir el concepto de \emph{ideología de la
lengua estándar} y vincularlo con la tradición sociolingüística clásica sobre
prestigio y estigma. Añade referencias canónicas de sociolingüística variacionista
y, si es posible, trabajos sobre el español.]

\subsection{Del juicio humano al sesgo algorítmico}
\label{sec:puente}

El argumento que articula esta tesis puede enunciarse como un puente en tres
pasos. Primero, existe una discriminación lingüística humana documentada que
sanciona ciertos sociolectos (Sección~\ref{sec:discriminacion-ling}). Segundo,
los sistemas de \PLN aprenden de texto producido y anotado por humanos, de modo
que pueden internalizar esas sanciones (Sección~\ref{sec:fuentes-sesgo}).
Tercero, existe evidencia empírica de que esto efectivamente ocurre: el ya
mencionado estudio de Sap y colaboradores \cite{Sap2019Risk} muestra un caso en
el que una variedad lingüística asociada a un grupo social recibe un trato
algorítmico sistemáticamente desfavorable. La hipótesis de trabajo de la tesis
es que un mecanismo análogo podría operar para los sociolectos del español
mexicano, hipótesis que el diseño experimental de los
Capítulos~\ref{cap:metodologia} y~\ref{cap:resultados} pone a prueba de forma
controlada.

\section{Clasismo estructural y lenguaje en el contexto mexicano}
\label{sec:clasismo}

\subsection{Clasismo estructural en México}
\label{sec:clasismo-mx}

El clasismo no se reduce a actitudes individuales de desprecio: es
\emph{estructural} en la medida en que se inscribe en instituciones, prácticas y
representaciones que distribuyen de manera desigual el reconocimiento y las
oportunidades según la posición social percibida. En el caso mexicano, el
clasismo se entrecruza además con marcadores raciales y geográficos, de modo que
clase, color de piel y origen tienden a leerse conjuntamente.

[FUENTE MX: desarrolla esta subsección con literatura mexicana específica
---estudios sobre discriminación (p.\ ej., informes de instituciones nacionales
dedicadas a prevenir la discriminación), trabajos sociológicos y antropológicos
sobre clasismo y colorismo. Esta subsección es el núcleo de tu objetivo
específico O1 y requiere fuentes que aún no están en referencias.bib. No incluí
citas aquí para no inventarlas.]

\subsection{Marcadores sociolingüísticos del clasismo en el español mexicano}
\label{sec:marcadores-mx}

En el español mexicano, ciertos rasgos lingüísticos funcionan como índices de
clase y son objeto de juicio social. Entre ellos suelen señalarse el contraste
entre léxico formal y coloquial, el uso de anglicismos como marca de cierto
capital cultural, fenómenos de pronunciación, marcadores discursivos
característicos y, en el registro escrito informal, la variación ortográfica.
Los términos \emph{naco} y \emph{fresa}, ampliamente usados en el habla
cotidiana, condensan precisamente la asociación entre formas de hablar (y de
presentarse) y posición social percibida, operando como etiquetas
estigmatizantes y de prestigio respectivamente.

Para los fines de esta tesis, la importancia de estos marcadores es metodológica:
constituyen las dimensiones de variación que el \emph{benchmark} controlado del
Capítulo~\ref{cap:metodologia} deberá operacionalizar, manteniendo constante el
contenido informativo. La selección final de los marcadores se justificará
empíricamente a partir de la encuesta de percepciones.

[FUENTE MX: respalda la caracterización de estos marcadores con trabajos
lingüísticos sobre el español de México; evita apoyarte solo en el conocimiento
común.]

\subsection{Implicaciones para sistemas automatizados desplegados en México}
\label{sec:implicaciones-mx}

Si los marcadores anteriores estructuran el juicio social humano y los sistemas
de \PLN aprenden de texto humano, cabe esperar que un clasificador entrenado
o ajustado con datos del español mexicano herede, al menos en parte, esas
asociaciones. El riesgo es concreto: en procesos como el filtrado automatizado
de currículums, una respuesta sistemáticamente distinta del modelo ante el
mismo contenido expresado en un registro estigmatizado se traduciría en
desventajas materiales para personas cuyo estilo de escritura se aparta del
registro dominante. Documentar y medir ese riesgo es el propósito de los
capítulos siguientes.

\section*{Síntesis del capítulo}
\addcontentsline{toc}{section}{Síntesis del capítulo}

El recorrido de este capítulo puede condensarse en cuatro proposiciones que
sostienen el resto de la tesis. Primera: el sesgo que interesa aquí es el
social, no el estadístico, y se vuelve evaluable cuando se fija un atributo
sensible y una métrica de comportamiento (Definición~\ref{def:sesgo}). Segunda:
ese sesgo puede entrar a un sistema de \PLN por cinco vías documentadas
---datos, anotaciones, representaciones, modelo y agenda de investigación---,
todas plausibles para el caso sociolectal. Tercera: la variación sociolectal es
variación de forma y no de contenido, de modo que un trato algorítmico
diferenciado ante ella no admite justificación informativa. Cuarta: en el
contexto mexicano existen marcadores lingüísticos de clase socialmente
sancionados, lo que hace verosímil ---aunque no demostrado--- que los sistemas
desplegados en ese contexto los hereden. Establecer si esto ocurre, y con qué
magnitud, es una pregunta empírica; el siguiente capítulo describe el diseño
con el que se aborda.
